% 
\chapter{Conclusões} 
\label{chap:Conc} % For referencing the chapter elsewhere, use Chapter~\ref{Conc}

O capítulo final, sintetiza as principais conclusões que se obteve neste projeto. Serão analisados os objetivos do capítulo ~\ref{chap:Introdução}.

%-------------------------------------------------------------------------------

\section{Objetivos concretizados} 
\label{sec:desc} %For referencing this section elsewhere, use Section~\ref{sec:desc}

Na Tabela~\ref{tab:grau} encontram-se os objetivos e o grau de conclusão dos mesmos.

\begin{table}[H]
\caption{Grau de Conclusão dos Objetivos}
\label{tab:grau}
\begin{center}
\begin{tabular}{|p{8cm}|c|}
\hline
\textbf{Objetivo} & \textbf{Conclusão} \\ \hline
Modelação e criação de escalões de acesso para controlar o acesso à transmissão & Concluído \\ \hline
Criação de canal único por criador & Concluído \\ \hline
Desenvolvimento da interface do criador & Maioritariamente Concluído \\ \hline
Desenvolvimento da interface do utilizador & Maioritariamente Concluído \\ \hline
Garantia de privacidade do Criador & Concluído \\ \hline
Implementação do \textit{chat} em direto & Concluído \\ \hline
\end{tabular}
\end{center}
\end{table}

\section{Limitações e trabalho futuro} 

Mesmo que a maior parte dos objetivos tenham sido concretizados, persistem limitações na experiência “ao vivo” relacionadas com latência e consistência de reprodução. A solução baseada em HLS de baixa latência é adequada para difusão massiva, mas não para interações sub-segundo. Recomenda-se a normalização do uso do \textit{IVS Player} e a reserva de \textit{hls.js} apenas como \textit{fallback} controlado. Propõe-se a criação de \textit{dashboards} operacionais com métricas de \textit{player} e rede do utilizador.

No \textit{chat} em tempo real, a arquitetura com API Gateway WebSocket e DynamoDB cumpre os requisitos, mas pode ser otimizada. Sugere-se \textit{heartbeats} periódicos para detetar ligações obsoletas e limitação de taxa por canal para mitigar \textit{backpressure}. A moderação pode evoluir de regras básicas para muting, \textit{rate limits} adaptativos e filtragem de \textit{spam}.

\section{Apreciação final} 

Em termos globais, considera-se que os objetivos principais do MVP são atingidos. O objetivo principal foi alcançado: realizar uma transmição em direto com poucos cliques. As restantes funcionalidades como os \textit{tiers}, o \textit{chat} e o menu de opções operam com tempos de resposta adequados e fornecem uma base sólida para evolução. 

O estágio permitiu ajustar e melhorar as boas práticas que desenvolvi ao longo do curso (DDD, CI/CD). Para além disso proporcionou uma experiencial de trabalho com a AWS, um serviço de computação na \textit{cloud}. Os serviços da AWS estendem-se para diversos ramos, \textit{data}, \textit{DevOps}, cibersegurança...

Por fim, a participação neste projeto de uma \textit{start-up}, trouxe momentos 

Por fim, a convivência direta na construção de uma \textit{start-up} proporcionou ganhos difíceis de reproduzir em contextos mais estáveis: proximidade real com a definição do produto e das prioridades, ciclos de iteração curtos orientados ao valor, gestão consciente dos compromissos entre rapidez e qualidade e responsabilidade ponta-a-ponta: do desenho ao \textit{deploy} e à monitorização. Essas aprendizagens, técnicas e organizacionais, sustentam a evolução futura do sistema e da equipa. 